\documentclass[12pt, a4paper]{article}

% Pakiety
\usepackage[utf8]{inputenc}
\usepackage[T1]{fontenc}
\usepackage[polish]{babel}
\usepackage{amsmath, amssymb}
\usepackage{graphicx}
\usepackage{float}
\usepackage{geometry}
\usepackage{hyperref}
\usepackage{booktabs}
\usepackage{caption}
\usepackage{listings}
\usepackage{xcolor}
\usepackage{parskip}

\geometry{
  a4paper,
  left=2.5cm,
  right=2.5cm,
  top=2.5cm,
  bottom=2.5cm
}

% Zakaz dzielenia wyrazów
\hyphenpenalty=10000
\exhyphenpenalty=10000
\emergencystretch=3em

\hypersetup{
  colorlinks=true,
  linkcolor=blue,
  urlcolor=blue,
  citecolor=blue
}

% Zmiana nazwy float figure: Rysunek -> Wykres
\addto\captionspolish{\renewcommand{\figurename}{Wykres}}

% Styl listingów
\lstset{
  language=Python,
  basicstyle=\ttfamily\small,
  keywordstyle=\color{blue}\bfseries,
  commentstyle=\color{gray}\itshape,
  stringstyle=\color{orange!80!black},
  numbers=left,
  numberstyle=\tiny\color{gray},
  numbersep=8pt,
  frame=single,
  breaklines=true,
  breakatwhitespace=true,
  captionpos=b,
  showstringspaces=false,
  columns=flexible,
  keepspaces=true,
  inputencoding=utf8,
  extendedchars=true,
  literate={ą}{{\k{a}}}1 {ę}{{\k{e}}}1 {ó}{{\'{o}}}1
            {ś}{{\'{s}}}1 {ł}{{\l{}}}1 {ź}{{\'{z}}}1
            {ż}{{\.z}}1 {ć}{{\'{c}}}1 {ń}{{\'{n}}}1
            {Ą}{{\k{A}}}1 {Ę}{{\k{E}}}1 {Ó}{{\'{O}}}1
            {Ś}{{\'{S}}}1 {Ł}{{L}}1,
}

% Strona tytułowa
\title{
  \textbf{Metody Obliczeniowe w Nauce i Technice} \\[0.5em]
  \large Laboratorium 7 \textendash{} Kwadratury, część II
}
\author{Mateusz Stoch}
\date{28 kwietnia 2026}

\begin{document}

\maketitle

% ============================================================
\section{Zadanie 1 \textendash{} Prędkość strzały z łuku}
% ============================================================

\subsection{Opis zadania}

Dane są wartości siły naciągu $F$ łuku w funkcji długości naciągu $x$:

\begin{center}
\begin{tabular}{r|rrrrrrrrrrr}
\toprule
$x$ [m] & 0.00 & 0.05 & 0.10 & 0.15 & 0.20 & 0.25 & 0.30 & 0.35 & 0.40 & 0.45 & 0.50 \\
\midrule
$F$ [N]  & 0    & 37   & 71   & 104  & 134  & 161  & 185  & 207  & 225  & 239  & 250  \\
\bottomrule
\end{tabular}
\end{center}

Łuk jest naciągnięty na odległość $L = 0{,}5\,\text{m}$, a masa strzały wynosi $m = 0{,}075\,\text{kg}$. Energia kinetyczna strzały w chwili wypuszczenia równa się pracy włożonej w naciągnięcie łuku:

\[
  \frac{1}{2}mv^2 = \int_0^L F\,dx
\]

\subsection{Argumentacja rozwiązania}

Dane są podane w postaci dyskretnych punktów pomiarowych, a nie analitycznej funkcji, dlatego całkę wyznaczono metodą Simpsona z biblioteki SciPy (\texttt{scipy.integrate.simpson}). Metoda Simpsona przy równoodległych węzłach jest odpowiednim wyborem, ponieważ minimalizuje błąd przy stosunkowo niewielu punktach pomiarowych.

\begin{lstlisting}[caption={Obliczenie pracy i prędkości strzały}]
x_data = np.array([0.00, 0.05, 0.10, ..., 0.50])
F_data = np.array([0, 37, 71, ..., 250])

W_simpson = np.float64(integrate.simpson(F_data, x=x_data))
v_simpson = np.sqrt(np.float64(2.0) * W_simpson / m)
\end{lstlisting}

\subsection{Wyniki}

\begin{figure}[H]
  \centering
  \includegraphics[width=0.85\textwidth]{wykres1.png}
  \caption{Siła naciągu łuku w zależności od długości naciągu. Zacieniowany obszar odpowiada pracy $W$.}
  \label{fig:luk}
\end{figure}

Obliczone wartości:
\begin{itemize}
  \item Praca: $W = 74{,}5333\,\text{J}$
  \item Prędkość strzały: $v = 44{,}5820\,\text{m/s} \approx 160\,\text{km/h}$
\end{itemize}

\subsection{Wnioski}

Metoda Simpsona dobrze nadaje się do całkowania danych dyskretnych. 
% ============================================================
\section{Zadanie 2 \textendash{} Ciało na sprężynie z tarciem}
% ============================================================

\subsection{Część (a): Prędkość $v_0$ w punkcie $x = 0$}

\subsubsection{Opis zadania}

Ciało o masie $m = 0{,}8\,\text{kg}$ jest przymocowane do sprężyny o stałej sprężystości $k = 80\,\text{N/m}$ i długości swobodnej $l = 0{,}4\,\text{m}$. Współczynnik tarcia między ciałem a poziomym prętem wynosi $\mu = 0{,}3$. Przyspieszenie ciała jest równe $\ddot{x} = -f(x)$, gdzie:
\[
  f(x) = \mu g + \frac{k}{m}\left(\mu l + x\right)\left(1 - \frac{l}{\sqrt{l^2 + x^2}}\right)
\]

Prędkość ciała w punkcie $x = 0$, gdy zaczyna ono ruch od $x = l$, wyraża się wzorem:
\[
  v_0 = \sqrt{2\int_0^l f(x)\,dx}
\]

\subsubsection{Argumentacja rozwiązania}

Funkcja podcałkowa $f(x)$ jest gładka na przedziale $[0, l]$, więc można bezpośrednio zastosować adaptacyjną kwadraturę Gaussa-Kronroda (\texttt{scipy.integrate.quad}).

\begin{lstlisting}[caption={Obliczenie prędkości v0}]
def f2(x):
    return mu*g + (k/m2)*(mu*l2 + x)*(1.0 - l2/np.sqrt(l2**2 + x**2))

W2, err2 = integrate.quad(f2, 0.0, float(l2))
v0 = np.sqrt(2.0 * W2)
\end{lstlisting}

\subsubsection{Wyniki}

\begin{figure}[H]
  \centering
  \includegraphics[width=0.85\textwidth]{wykres2.png}
  \caption{Funkcja przyspieszenia $f(x)$ na przedziale $[0, l]$.}
  \label{fig:sprezyna_a}
\end{figure}

\begin{itemize}
  \item Całka: $\int_0^l f(x)\,dx = 3{,}119190$
  \item Szacowany błąd: $3{,}46 \times 10^{-14}$
  \item Prędkość: $v_0 = 2{,}497675\,\text{m/s}$
\end{itemize}

\subsection{Część (b): Stała C i czas $t$}

\subsubsection{Opis zadania}

Czas potrzebny na przebycie drogi od $x = l$ do $x = 0$ wynosi $t = C\sqrt{m/k}$, gdzie:
\[
  C = \int_0^1 \left[(\sqrt{2}-1)^2 - (\sqrt{1+z}-1)^2\right]^{-1/2} dz
\]

Funkcja podcałkowa ma osobliwość w punkcie $z = 1$, która zachowuje się jak $(1-z)^{-1/2}$.

\subsubsection{Argumentacja rozwiązania}

Aby poradzić sobie z osobliwością, zastosowano następujące podejście:
\begin{enumerate}
  \item Bezpośrednie wywołanie \texttt{integrate.quad} z parametrem \texttt{points=[0.99]} -- wymuszenie podziału siatki w pobliżu osobliwości.
\end{enumerate}

\begin{lstlisting}[caption={Obliczenie stałej C z obsługą osobliwości}]
def integrand_C(z):
    A = (np.sqrt(2.0) - 1.0)**2
    B = (np.sqrt(1.0 + z) - 1.0)**2
    val = np.where(A - B <= 0.0, 1e-30, A - B)
    return 1.0 / np.sqrt(val)

C_val, err_C = integrate.quad(
    integrand_C, 0.0, 1.0,
    limit=200, points=[0.99],
    epsabs=1e-10, epsrel=1e-10
)
t_val = C_val * np.sqrt(m2 / k)
\end{lstlisting}

\subsubsection{Wyniki}

\begin{itemize}
  \item Stała $C = 3{,}970020$
  \item Szacowany błąd: $1{,}24 \times 10^{-10}$
  \item Czas: $t = C\sqrt{m/k} = 3{,}970020 \cdot 0{,}100000 = 0{,}397002\,\text{s}$
\end{itemize}

\subsection{Wnioski}

Zadanie 2 pokazuje dwa typowe problemy praktyczne przy numerycznym całkowaniu:
\begin{itemize}
  \item W części (a) funkcja jest gładka, co sprawia, że adaptacyjna kwadratura Gaussa-Kronroda daje wynik niemal natychmiastowo z precyzją rzędu błędu maszynowego.
  \item W części (b) istnieje osobliwość w końcowym punkcie przedziału, która wymaga specjalnego traktowania. Poinformowanie algorytmu o lokalizacji osobliwości (parametr \texttt{points}).
\end{itemize}

% ============================================================
\section{Zadanie 3 \textendash{} Wahadło matematyczne}
% ============================================================

\subsection{Opis zadania}

Dane jest wahadło matematyczne o długości $l = 1\,\text{m}$ i wychyleniu początkowym $\theta_0$. Dla małych wychyleń okres jest przybliżany wzorem:
\[
  T_0 = 2\pi\sqrt{\frac{l}{g}}
\]

Dla dowolnych wychyleń dokładny okres wyraża się przez całkę eliptyczną pierwszego rodzaju $K(k)$:
\[
  T(\theta_0) = 4\sqrt{\frac{l}{g}}\,K(k), \qquad K(k) = \int_0^{\pi/2}\frac{d\theta}{\sqrt{1 - k^2\sin^2\theta}}, \qquad k = \sin\frac{\theta_0}{2}
\]

\subsection{Część (a): Całka eliptyczna kwadraturą}

\subsubsection{Argumentacja rozwiązania}

Całkę eliptyczną wyznaczono przy użyciu \texttt{scipy.integrate.quad} z funkcją podcałkową $\frac{1}{\sqrt{1-k^2\sin^2\theta}}$ na przedziale $[0, \pi/2]$. Jako wartość referencyjną przyjęto wartość zwracaną przez funkcję \texttt{scipy.special.ellipk}.

\begin{lstlisting}[caption={Obliczenie całki eliptycznej kwadraturą}]
k_mod = np.sin(theta0 / 2.0)
K_true = np.float64(special.ellipk(k_mod**2))

def elliptic_integrand(theta, k):
    return 1.0 / np.sqrt(1.0 - k**2 * np.sin(theta)**2)

K_quad, err_K = integrate.quad(
    elliptic_integrand, 0.0, np.pi/2, args=(float(k_mod),)
)
T_true_quad = 4.0 * np.sqrt(l3 / g3) * K_quad
\end{lstlisting}

\subsubsection{Wyniki}

Dla $\theta_0 = 45^{\circ}$:
\begin{itemize}
  \item Moduł: $k = \sin(22{,}5^{\circ}) \approx 0{,}38268343$
  \item $K(k)$ z \texttt{scipy.special.ellipk}: $1{,}63358631$
  \item $K(k)$ kwadratura: $1{,}63358631$
  \item Błąd względny (kwadratura vs scipy): $1{,}36 \times 10^{-16}$
  \item Dokładny okres: $T(\theta_0 = 45^{\circ}) = 2{,}086256\,\text{s}$
  \item Przybliżony okres (małe drgania): $T_0 = 2{,}006067\,\text{s}$
\end{itemize}

\subsection{Część (b): Całka eliptyczna przez AGM}

\subsubsection{Opis zadania}

Istnieje wzór na całkę eliptyczną przez średnią arytmetyczno-geometryczną (AGM):
\[
  K(k) = \frac{\pi}{2} \cdot \frac{1}{\mathrm{AGM}(1+k,\, 1-k)}
\]

Ciągi $a_n$, $b_n$ definiujące AGM spełniają rekurencję:
\[
  a_0 = 1+k, \quad b_0 = 1-k, \qquad a_{n+1} = \frac{a_n+b_n}{2}, \quad b_{n+1} = \sqrt{a_n b_n}
\]

\subsubsection{Argumentacja rozwiązania}

AGM zbiega kwadratowo, więc kilka iteracji wystarczy, by osiągnąć precyzję maszynową. Zbieżność jest jeszcze szybsza niż wykładnicza.

\begin{lstlisting}[caption={Obliczenie K(k) przez AGM}]
def agm_K(k_val, n_iter):
    a = 1.0 + k_val
    b = 1.0 - k_val
    for _ in range(n_iter):
        a_new = (a + b) / 2.0
        b_new = np.sqrt(a * b)
        a, b = a_new, b_new
    return (np.pi / 2.0) / a
\end{lstlisting}

\subsubsection{Wyniki}

\begin{figure}[H]
  \centering
  \includegraphics[width=0.85\textwidth]{wykres3.png}
  \caption{Błąd względny przybliżenia $K(k)$ przez AGM w zależności od liczby iteracji ($\theta_0 = 45^{\circ}$).}
  \label{fig:agm}
\end{figure}

Po 20 iteracjach:
\begin{itemize}
  \item $K(k)_\mathrm{AGM} = 1{,}63358631$
  \item Błąd względny AGM: $0{,}00 \times 10^{0}$ (poniżej progu precyzji maszynowej)
  \item Okres AGM: $T = 2{,}086256\,\text{s}$
\end{itemize}

Z wykresu wynika, że AGM osiąga precyzję maszynową już po około 4--5 iteracjach.

\subsection{Błąd przybliżenia małych drgań}

\begin{figure}[H]
  \centering
  \includegraphics[width=0.85\textwidth]{wykres4.png}
  \caption{Błąd względny okresu $T_0$ (przybliżenie małych drgań) w zależności od amplitudy $\theta_0$.}
  \label{fig:blad_t0}
\end{figure}

Wykres pokazuje, że błąd wzoru $T_0 = 2\pi\sqrt{l/g}$ rośnie gwałtownie wraz ze wzrostem kąta wychylenia. Dla małych kątów (do ok. $15^{\circ}$) błąd jest pomijalny, natomiast dla $\theta_0 = 90^{\circ}$ przekracza już kilkanaście procent.

\subsection{Wnioski}

\begin{enumerate}
  \item Kwadratura adaptacyjna Gaussa-Kronroda (\texttt{quad}) wyznacza całkę eliptyczną z precyzją bliską granicy maszynowej.
  \item Metoda AGM zbiega znacznie szybciej niż typowe kwadratury numeryczne, ponieważ jej zbieżność jest kwadratowa. W praktyce 5 iteracji wystarczy do osiągnięcia pełnej precyzji \texttt{float64}.
  \item Przybliżenie okresu małych drgań $T_0$ jest wystarczająco dobre tylko dla kątów wychylenia nieprzekraczających kilkunastu stopni. Przy większych kątach należy korzystać ze wzoru dokładnego opartego na całce eliptycznej.
\end{enumerate}

% ============================================================
\section{Podsumowanie}
% ============================================================

Laboratorium 7 skupiło się na zastosowaniach całkowania numerycznego do problemów fizycznych. Wykazano, że:
\begin{itemize}
  \item Do całkowania danych tabelarycznych (zadanie 1) wygodnie stosuje się prostą kwadraturę Simpsona.
  \item Adaptacyjna kwadratura Gaussa-Kronroda (\texttt{scipy.integrate.quad}) radzi sobie znakomicie z gładkimi funkcjami i funkcjami z osobliwościami po wskazaniu ich lokalizacji (zadanie 2).
  \item Dla specjalnych funkcji matematycznych (całka eliptyczna) istnieją efektywne algorytmy iteracyjne (AGM), które osiągają bardzo wysoką precyzję przy minimalnej liczbie obliczeń (zadanie 3).
\end{itemize}

\end{document}
